% !TEX TS-program = pdflatex
% !TEX encoding = UTF-8 Unicode
% Enhanced-dissipation tools: a page of the wiki "Enhanced dissipation for passive scalars", generated from the wiki on 12 November 2026.
% It compiles as it is with pdflatex (Overleaf, TeXShop, TeXstudio, ...); typeset twice for the citations.
% To propose a change: keep the preamble and the line \title{...}; keep only the blocks you change
% (a \section{...} with its content, or a statement or proof with its \label{...}), change them, and post
% the file in a new topic of the Zulip channel of the part. The rules: page "How to contribute" of the wiki.
\documentclass[11pt]{article}
\usepackage[utf8]{inputenc}
\usepackage[T1]{fontenc}
\usepackage{lmodern}
\usepackage[a4paper,margin=2.5cm]{geometry}
\usepackage{amsmath,amssymb,amsthm}
\usepackage{graphicx,tabularx,listings}
\usepackage[skins,breakable]{tcolorbox}
\usepackage[round]{natbib}
\usepackage[hidelinks]{hyperref}
\setlength{\emergencystretch}{2em}

\newtheorem{theorem}{Theorem}
\newtheorem{proposition}[theorem]{Proposition}
\newtheorem{lemma}[theorem]{Lemma}
\newtheorem{corollary}[theorem]{Corollary}
\newtheorem{conjecture}[theorem]{Conjecture}
\newtheorem*{claim}{Claim}
\theoremstyle{definition}
\newtheorem{definition}[theorem]{Definition}
\newtheorem{example}[theorem]{Example}
\theoremstyle{remark}
\newtheorem{remark}[theorem]{Remark}

%% The commands of the wiki. This preamble is the same for every page: do not change it, and do not
%% define commands of your own (the wiki has none).
\makeatletter
% after a statement: \begin{idea}, \begin{architecture}, \begin{note}[Title] (and \begin{proof}[Title])
\newcommand{\ed@runin}[1]{\par\topsep6\p@\@plus6\p@\relax\trivlist\item[\hskip\labelsep\itshape#1\@addpunct{.}]\ignorespaces}
\newenvironment{idea}[1][Idea]{\ed@runin{#1}}{\endtrivlist\@endpefalse}
\newenvironment{architecture}[1][Proof architecture]{\ed@runin{#1}}{\endtrivlist\@endpefalse}
% \begin{note}[Title]: an aside, in a dashed box; its title in italics at the start of the text
\tcbset{ed@note/.style={enhanced jigsaw,breakable,frame hidden,interior hidden,arc=2pt,borderline={0.5pt}{0pt}{black!50,dash pattern=on 3pt off 2pt},
  left=3mm,right=3mm,top=2mm,bottom=2mm,before skip=3mm plus 1mm,after skip=3mm plus 1mm}}
\newenvironment{note}[1][Note]{\begin{tcolorbox}[ed@note]\noindent{\itshape#1\@addpunct{.}}\enspace\ignorespaces}{\end{tcolorbox}}
% text still to be written; labels of the text; the Lean banner of the statement above
\newcommand{\placeholder}[1]{\if@inlabel\else\par\noindent\fi\textit{[#1]}\@ifnextchar\end{}{\par}}
\newcommand{\ed@label}[1]{\mbox{\footnotesize[#1]}}
\newcommand{\demo}[1]{\ed@label{demo text}}
\newcommand{\heuristic}[1]{\ed@label{heuristic}}
\newcommand{\numerics}[1]{\ed@label{numerics}}
\newcommand{\published}[1]{\ed@label{published}}
\newcommand{\preprint}[1]{\ed@label{preprint, to be verified}}
\newcommand{\raw}[1]{\ed@label{raw, unchecked}}
\def\ed@formalized{formalized}\def\ed@image{image}\def\ed@video{video}\def\ed@diagram{diagram}
\newcommand{\lean}[2][]{\par{\noindent\footnotesize\def\ed@a{#2}Lean: \ifx\ed@a\ed@formalized formalized\ifx\relax#1\relax\else\
  (\href{#1}{certificate})\fi\else not yet formalized\fi.\par}}
% the contributors of a section, most recent first (chosen by the editors), in small type under its title
\newcommand{\contributors}[1]{\par{\noindent\footnotesize\itshape Contributors: #1\par}\smallskip}
% the two lines under the title: the page's team and status, its Zulip topic (default: the title)
\def\ed@title{}\let\ed@maketitle\maketitle\renewcommand{\maketitle}{\let\ed@title\@title\ed@maketitle}
\newcommand{\pagestatus}[1]{{\par\centering\small#1\par}}
\newcommand{\zulip}[2][]{{\par\centering\small Discuss on Zulip: \##2, topic
  ``\ifx\relax#1\relax\ed@title\else#1\fi''.\par}\medskip}
% links to other pages of the wiki, anchors, wiki code shown as text
\newcommand{\wikilink}[2]{\ifx\relax#2\relax#1\else#2\fi}
\newcommand{\wikianchor}[1]{\phantomsection\label{#1}}
\newcommand{\wikipage}[1]{}% the first line of a page file of the whole-wiki project (pages/<page>.tex): the page's title
\newcommand{\nowikicode}[1]{\texttt{\frenchspacing#1}}
\newcommand{\nowiki}[1]{#1}
\newcommand{\wikitemplate}[1]{\texttt{\{\{#1\}\}}}
\newcommand{\wikidiv}[2]{\par#2\par}
% \begin{quote}: a sentence set apart, such as the central question of the Overview; centred, in italics
\renewenvironment{quote}{\list{}{\rightmargin\leftmargin}\item\relax\centering\itshape}{\endlist}
% pictures, videos, diagrams: a grey panel with the picture and its caption. The picture (for a video, its poster image)
% is shown when its file, or a .pdf or .png file of the same name, is next to this file, in figures/ or in ../media
\tcbset{ed@media/.style={enhanced,colback=black!4,frame hidden,arc=2pt,left=3mm,right=3mm,top=2.5mm,bottom=2mm,
  before skip=3mm plus 1mm,after skip=3mm plus 1mm}}
\def\ed@poster{}\def\ed@record{}\def\ed@svg{svg}
\define@key{ed@media}{poster}{\def\ed@poster{#1}}\define@key{ed@media}{record}{\def\ed@record{#1}}\define@key{ed@media}{id}{}
\newcommand{\ed@find}[1]{\IfFileExists{#1}{\edef\ed@found{#1}}{\IfFileExists{figures/#1}{\edef\ed@found{figures/#1}}{%
  \IfFileExists{../media/#1}{\edef\ed@found{../media/#1}}{}}}}
\newcommand{\ed@picturefile}[1]{\def\ed@found{}\edef\ed@b{#1}\ifx\ed@b\@empty\else\filename@parse{\ed@b}\edef\ed@base{\filename@base}%
  \edef\ed@ext{\filename@ext}\ifx\ed@ext\ed@svg\else\ed@find{\ed@b}\fi
  \ifx\ed@found\@empty\ed@find{\ed@base.pdf}\fi\ifx\ed@found\@empty\ed@find{\ed@base.png}\fi\fi}
\newcommand{\ed@panel}[2]{\begin{tcolorbox}[ed@media,#1]\ifx\ed@found\@empty\else{\centering\includegraphics[width=\linewidth,
  height=0.35\textheight,keepaspectratio]{\ed@found}\par}\smallskip\fi{\small\textbf{\ifx\ed@a\ed@video Video\else\ifx\ed@a\ed@diagram
  Diagram\else Figure\fi\fi.}\enspace#2\par}{\footnotesize\itshape\ifx\ed@a\ed@video The video plays in the online edition of this page.\fi
  \ifx\ed@record\@empty\else\ifx\ed@a\ed@video\ \fi How it was made: page Numerical methods, record \texttt{\ed@record}.\fi\par}
  \end{tcolorbox}}
% in the text, the panel is a figure (it floats to where it fits); inside a box (a note, a table), it stays in place
\renewcommand{\topfraction}{0.85}\renewcommand{\bottomfraction}{0.6}\renewcommand{\textfraction}{0.1}\renewcommand{\floatpagefraction}{0.8}
\newcommand{\media}[4][]{\par\def\ed@poster{}\def\ed@record{}\setkeys{ed@media}{#1}\def\ed@a{#2}%
  \ifx\ed@a\ed@video\ed@picturefile{\ed@poster}\else\ed@picturefile{#3}\fi
  \ifinner\ed@panel{}{#4}\else\begin{figure}[htbp]\ed@panel{before skip=0pt,after skip=0pt}{#4}\end{figure}\fi}
% tables (\lv{label}{text}; \header{a & b}, \row{a & b}) and code
\newenvironment{levels}[1][]{\par\smallskip\noindent\tabularx{\linewidth}{@{}>{\bfseries\raggedright\arraybackslash}p{0.2\linewidth}X@{}}}{\endtabularx\par\smallskip}
\newcommand{\lv}[2]{#1 & #2\\[3pt]}
\newenvironment{wikitable}[1]{\par\smallskip\noindent\tabularx{\linewidth}{|*{#1}{>{\raggedright\arraybackslash}X|}}\hline}{\endtabularx\par\smallskip}
\newcommand{\header}[1]{#1\\\hline\hline}
\newcommand{\row}[1]{#1\\\hline}
\lstnewenvironment{wikicode}[1][]{\lstset{basicstyle=\ttfamily\footnotesize,breaklines=true,columns=fullflexible,
  keepspaces=true,showstringspaces=false,frame=single}\ifx\relax#1\relax\else\lstset{language=#1}\fi}{}
% the heading of the reference list is the section References of the page
\renewcommand{\bibsection}{}
% characters of the wiki's text that LaTeX does not know
\DeclareUnicodeCharacter{2192}{\ensuremath{\rightarrow}}\DeclareUnicodeCharacter{203A}{\ensuremath{\rangle}}
\DeclareUnicodeCharacter{2713}{\ensuremath{\checkmark}}\DeclareUnicodeCharacter{25CB}{\ensuremath{\circ}}
\DeclareUnicodeCharacter{2248}{\ensuremath{\approx}}\DeclareUnicodeCharacter{2264}{\ensuremath{\leq}}
\DeclareUnicodeCharacter{2265}{\ensuremath{\geq}}\DeclareUnicodeCharacter{2212}{\ensuremath{-}}
\makeatother

\begin{document}
\title{Enhanced-dissipation tools}
\author{Wiki \emph{Enhanced dissipation for passive scalars}}
\date{12 November 2026}
\maketitle
\pagestatus{Contributors: to be named · Status: skeleton}
\zulip{toolkit}

Tools specific to enhanced dissipation. A statement comes to the Toolkit when several proofs use it, or when it holds in more generality than one result; a step used by one proof only is a claim inside that proof. Each statement is given once, with its proof or a reference, and the proofs elsewhere link to it.

\section{Tools}
\placeholder{For instance spectral-gap and resolvent lemmas, and the link between mixing and enhanced dissipation.}

\subsection{From resolvent bounds to semigroup bounds}
\begin{definition}[m-accretive operator]\label{def-m-accretive}
A closed operator $H$ on a complex Hilbert space $X$ is m-accretive if, for every $\mu\in\mathbb{C}$ with $\operatorname{Re}\mu>0$, the operator $H+\mu$ is invertible and $\|(H+\mu)^{-1}\|\leq(\operatorname{Re}\mu)^{-1}$. Then $D(H)$ is dense, $\operatorname{Re}\langle Hf,f\rangle\geq0$ for $f\in D(H)$, and $-H$ generates a semigroup of contractions $e^{-tH}$.
\end{definition}
\begin{theorem}[Gearhart–Prüss type bound, Wei]\label{thm-gearhart-pruss-wei}
Let $H$ be an m-accretive operator on $X$, and let
\[\Psi(H):=\inf\big\{\|(H-i\lambda)f\|\ :\ f\in D(H),\ \|f\|=1,\ \lambda\in\mathbb{R}\big\}.\]
Then $\|e^{-tH}\|\leq e^{\pi/2}\,e^{-\Psi(H)t}$ for all $t\geq0$.
\end{theorem}
\begin{idea}
Cut the solution in time: $\chi(s)e^{-sH}f$ solves the equation with source $\chi'(s)e^{-sH}f$. After a Fourier transform in time the resolvent on the imaginary axis appears, and by Plancherel the definition of $\Psi$ becomes a weighted inequality for the nonincreasing function $g(s)=\|e^{-sH}f\|^2$. With $\chi$ exponential in the bulk this forces $g$ to decay like $e^{-2\Psi s}$. Bringing $\chi$ to zero at the two ends costs the constant $e^{\pi/2}$.
\end{idea}
\begin{proof}\label{proof-gearhart-pruss-wei}
\emph{Sources.} \citet{Wei2021Resolvent}, Section 2, partly motivated by \citet{HelfferSjostrand2010}.

Let $\Psi=\Psi(H)$. Since $D(H)$ is dense, it is enough to prove $\|e^{-tH}f\|\leq e^{\pi/2-\Psi t}\|f\|$ for $f\in D(H)$. Fix such an $f$ and let $g(s)=\|e^{-sH}f\|^2$, which is nonincreasing because the semigroup is a contraction. In particular there is nothing to prove if $\Psi t\leq\pi/2$, so we assume $\Psi t>\pi/2$.

\textbf{Step 1 (an inequality for $g$).}
\begin{claim}\label{clm-gearhart-pruss-wei}
For every $l>0$ and every real $\chi\in C^1([0,l])$ with $\chi(0)=\chi(l)=0$,
\[\int_0^l\chi'(s)^2g(s)\,ds\geq\Psi^2\int_0^l\chi(s)^2g(s)\,ds.\]
\end{claim}
\begin{proof}[Proof of the claim]\label{proof-clm-gearhart-pruss-wei}
Let $u(s)=\chi(s)e^{-sH}f$ and $w(s)=\chi'(s)e^{-sH}f$ on $[0,l]$, so that $\partial_su+Hu=w$. For $\lambda\in\mathbb{R}$ let $\hat u(\lambda)=\int_0^lu(s)e^{-i\lambda s}\,ds$, and similarly $\hat w$. Since $H$ is closed, $\hat u(\lambda)\in D(H)$. Integrating by parts in $s$, the boundary terms vanish because $\chi(0)=\chi(l)=0$, and
\[(H+i\lambda)\hat u(\lambda)=\hat w(\lambda).\]
By the definition of $\Psi$, $\|\hat w(\lambda)\|\geq\Psi\|\hat u(\lambda)\|$ for every $\lambda$. By Plancherel's theorem for functions with values in $X$,
\[\int_0^l\|w(s)\|^2\,ds=\frac{1}{2\pi}\int_{\mathbb{R}}\|\hat w(\lambda)\|^2\,d\lambda\geq\frac{\Psi^2}{2\pi}\int_{\mathbb{R}}\|\hat u(\lambda)\|^2\,d\lambda=\Psi^2\int_0^l\|u(s)\|^2\,ds.\]
\end{proof}

\textbf{Step 2 (choice of $\chi$).} Let
\[t_1=\frac{\pi}{4\Psi},\qquad t_2=t-\frac{\pi}{4\Psi},\qquad t_3=t+\frac{\pi}{4\Psi},\qquad l=t+\frac{\pi}{2\Psi},\]
so that $t_1<t_2$, and
\[\chi(s)=\begin{cases}\sin(\Psi s),&0\leq s\leq t_1,\\ \frac{1}{\sqrt2}\,e^{\Psi s-\pi/4},&t_1\leq s\leq t_2,\\ e^{\Psi t-\pi/2}\cos(\Psi(s-t)),&t_2\leq s\leq l.\end{cases}\]
This function is $C^1$ and vanishes at $0$ and $l$. Setting $h:=\chi'^2-\Psi^2\chi^2$, the claim reads $\int_0^lh\,g\,ds\geq0$, and
\[h(s)=\begin{cases}\Psi^2\cos(2\Psi s),&0\leq s\leq t_1,\\ 0,&t_1\leq s\leq t_2,\\ -\Psi^2e^{2\Psi t-\pi}\cos(2\Psi(s-t)),&t_2\leq s\leq l.\end{cases}\]
Hence $h\geq0$ on $[0,t_1]$ and on $[t_3,l]$, $h\leq0$ on $[t_2,t_3]$, and
\[\int_0^{t_1}h\,ds=\frac{\Psi}{2},\qquad\int_{t_2}^{t}h\,ds=-\frac{\Psi}{2}e^{2\Psi t-\pi},\qquad\int_t^lh\,ds=0.\]

\textbf{Step 3 (conclusion).} Since $g$ is nonincreasing, $hg\leq h\,g(0)$ on $[0,t_1]$, $hg\leq h\,g(t)$ on $[t_2,t]$, and $hg\leq h\,g(t_3)$ on $[t,l]$. Therefore
\[0\leq\int_0^lh\,g\,ds\leq\frac{\Psi}{2}\,g(0)-\frac{\Psi}{2}\,e^{2\Psi t-\pi}g(t),\]
that is $g(t)\leq e^{\pi-2\Psi t}g(0)$, which is the bound.
\end{proof}


\section*{References}
%%wiki <!-- refs-begin: Wei2021Resolvent HelfferSjostrand2010 -->
\begin{thebibliography}{}
\bibitem[Helffer and Sjöstrand(2010)]{HelfferSjostrand2010} B. Helffer and J. Sjöstrand, \emph{From resolvent bounds to semigroup bounds}, preprint (2010), \href{https://arxiv.org/abs/1001.4171}{arXiv:1001.4171}.
\bibitem[Wei(2021)]{Wei2021Resolvent} D. Wei, \emph{Diffusion and mixing in fluid flow via the resolvent estimate}, Science China Mathematics \textbf{64} (2021), 507–518, \href{https://arxiv.org/abs/1811.11904}{arXiv:1811.11904}.
\end{thebibliography}
%%wiki <!-- refs-end -->
\end{document}
